\chapter{Testování aplikace}
Testování aplikace bylo zaměřeno na ověření správnosti implementace, stability aplikace a odhalení chyb vznikajících při běžném i méně obvyklém použití.

V průběhu vývoje byly využity různé přístupy k testování, od automatizovaných jednotkových testů přes manuální testování až po sledování chyb v reálném provozu pomocí nástrojů třetích stran.


\section{Jednotkové testy}
Pro ověření správnosti aplikační logiky byly využity jednotkové testy, zaměřené především na doménovou vrstvu, zejména use-case třídy.

Testy byly navrženy tak, aby ověřovaly chování jednotlivých use-case tříd nezávisle na konkrétní implementaci datových zdrojů.
Díky oddělení doménové a datové vrstvy bylo možné testovat use-case třídy pomocí mockovaných repozitářů, což umožnilo simulovat různé scénáře, včetně chybových stavů.

Největší část testů je zaměřena na funkční celky \texttt{core} a \texttt{shop},
které obsahují klíčové funkce aplikace.

Testy pomohly odhalit například chyby v mapování doménových chyb, kdy use-case třídy po některých úpravách vracely nepřesné nebo příliš obecné chyby (například obecnou chybu pro práci s obchody namísto chyby chybějící autentizace uživatele).


\section{Firebase crashlytics}
Pro sledování chyb v reálném provozu aplikace je využívána služba Firebase Crashlytics, která umožňuje automatické zaznamenávání pádů aplikace.

Crashlytics poskytuje detailní informace o vzniklých chybách, jako jsou informace o zařízení, verze aplikace nebo konkrétní zdroj chyby ve zdrojovém kódu, což umožňuje rychle identifikovat příčinu problému a následně jej opravit.


\section{Manuální testování}
% testovací scénáře -> nalezené chyby v konkrétním scénáři, které jsem opravil
% například:
%  - navigace - navigace zpět po přihlášení vedla na přihlašovací obrazovku - správný navigační graf
%  - reset mapy při přechodu na recenze obchodu
%  - přídání confirm dialogů


\section{Testování na uživatelích}
% scénáře
% na technicky zdalých i nezdalých uživatelích
% zjištěné problémy a opravy
% doplnění popisků, upozornění